\setscript[hanzi]
\setupalign[hanging, hz]

\usemodule[zhfontsext]
\usetypescript[cnfontb]
\setupbodyfont[cnfontb,rm,12pt] % 小四号 = 12pt
\mainlanguage[cn]

\setupinterlinespace[line=3.66ex] %line=3.05ex相当于WORD 中的1.0倍行距, 3.66ex相当于WORD 中的1.2倍行距

\setuppapersize[A4]
\setuplayout[width=fit,height=fit,backspace=10mm,cutspace=10mm,topspace=10mm,margin=10mm,header=10mm,footer=10mm]
\setuppagenumbering[location={header,right},style=\bfb]

\define[1]\exampleLine{
    \blank[5pt]
    \framed[frame=off,align=flushleft,offset=3pt,strut=none,width=\textwidth,height=fit,background=color,backgroundcolor=gray]{{\bfx #1}}
    \blank[5pt]
}

\starttext

\exampleLine{代码示例 - 5:}

\define[1]\colorEnum{\darkgreen{#1}}
\defineenumeration[question]
\setupenumeration[question][alternative=top,left=第,right=题,text=大题,title=yes,inbetween={\blank[none]},after=\blank]
\setupenumeration[subquestion][alternative=serried,width=fit,text=小题,color=darkblue,margin=10mm,headcommand=\colorEnum]

\startquestion[title={共十分}] 地理常识题.
\stopquestion

\startsubquestion
河北省会是？
\stopsubquestion
\startsubquestion
安徽省会是？
\stopsubquestion
\resetsubquestion
\startsubquestion
河南省会是？
\stopsubquestion 

\stoptext